% Exact restricted classification candidate; independent whole-proof audit requested.
\subsection{The twenty-six remaining ternary-glue cycles}
\label{subsec:k12-ternary-glue}

The proper-sublattice case in Proposition~\ref{prop:k12-trimodular-generated}
also admits a finite classification. The following statement identifies
actual cells and their boundaries. It does not assert that their homology
classes are nonzero.

\begin{theorem}[K17-K12-TERNARY-CYCLES]
\label{thm:k12-ternary-cycles}
Let $\rho$ be a rank-ten projective degree-twelve Voronoi cell whose
largest minimum-basis index is three and whose minimum vectors generate
a proper sublattice $L_0$ of its ambient lattice $L$.
Then either its orientation coinvariant is zero, or it is integrally
equivalent to exactly one of the twenty-six templates in
Table~\ref{tab:k12-ternary-templates} below. Each listed template is
realized by a positive form with exactly thirteen minimum pairs, has
orientation-preserving stabilizer, and is a cycle over
$R=\mathbb Z[1/13!]$.

Consequently, after subtracting the previously certified boundary
$\sigma$, every degree-twelve cycle supported on cells of maximal
minimum-basis index at most three is represented by an $R$-linear
combination of these twenty-six cycles. No linear independence or
nonvanishing of their homology classes is asserted.
\end{theorem}

\paragraph{Template notation.}
Write the following lists as columns:
\[
\begin{split}
\mathcal B_1&=((0,0,1),(0,1,0),(1,0,-1),(1,0,0),
                  (1,1,-1),(1,1,0)),\\
\mathcal B_2&=((0,0,1),(0,1,0),(0,1,1),(1,0,0),(1,0,1)),\\
\mathcal B_3&=((0,0,1),(0,1,0),(0,1,1),(1,0,0)),\\
\mathcal B_4&=((0,0,1),(0,1,0),(1,-1,-1),(1,0,0)),\\
\mathcal B_5&=((0,0,1),(0,1,0),(1,0,0)).
\end{split}                                                    \tag{TG.4}
\]
A mask $m\in\{1,\ldots,7\}$ denotes the subset
$C(m)=\{c\in\mathbb F_3:\text{bit }c\text{ of }m\text{ is one}\}$.
Thus $1,3,5,6,7$ denote $\{0\},\{0,1\},\{0,2\},\{1,2\}$
and $\mathbb F_3$, respectively.
For a mask list on $\mathcal B_i$, repeat its $j$th column once for
each $c\in C(m_j)$ and give that position color $c$.
The last table column specifies additional zero Gale columns:
$-$ means none, $0$ means one of color zero, $1$ means one of color
one, and $01$ means two of colors zero and one.
The resulting $3$-by-$13$ matrix is $B$, and the color vector is $c$.
Define
\[
 L_0=\mathbb Z^{13}/\operatorname{row}(B),\qquad
 L=L_0+\mathbb Z\frac{\bar c}{3}\ \subset L_0\otimes\mathbb Q.
                                                               \tag{TG.5}
\]
The thirteen images of the coordinate vectors give the template rays.
Every $B$ used here has a unit maximal minor, so its row lattice is
saturated and $L_0$ is free of rank ten.

\begin{table}[ht]
\centering
\begin{tabular}{ccl|ccl}
Template&Mask list&Zero colors&Template&Mask list&Zero colors\\\hline
$1.1$&$113777$&$-$&$1.11$&$131757$&$1$\\
$1.2$&$131377$&$1$&$1.12$&$131773$&$1$\\
$1.3$&$131637$&$01$&$1.13$&$131775$&$1$\\
$1.4$&$131657$&$01$&$1.14$&$131777$&$-$\\
$1.5$&$131673$&$01$&$1.15$&$133337$&$1$\\
$1.6$&$131675$&$01$&$1.16$&$133537$&$1$\\
$1.7$&$131677$&$0$&$1.17$&$133557$&$1$\\
$1.8$&$131677$&$1$&$1.18$&$133567$&$1$\\
$1.9$&$131737$&$1$&$1.19$&$133577$&$-$\\
$1.10$&$131756$&$01$&$1.20$&$173537$&$-$\\\hline
$2.1$&$31737$&$01$&$2.4$&$71733$&$01$\\
$2.2$&$33337$&$01$&$2.5$&$71737$&$0$\\
$2.3$&$71337$&$01$&$2.6$&$71737$&$1$
\end{tabular}
\caption{The surviving proper-index-three templates. The first index
specifies $\mathcal B_1$ or $\mathcal B_2$ in (TG.4).}
\label{tab:k12-ternary-templates}
\end{table}

\begin{proof}
The index $[L:L_0]$ divides every minimum-basis index and divides an
index equal to three. Since it is proper, it is three. Every independent
ten minima are consequently a basis of $L_0$. Choose one of these
bases and write the minimum matrix as $V=[I_{10}\ A]$.
Every square minor of $A$ is, up to sign, a maximal minor of $V$;
thus $V$ and its integral relation matrix $[-A^t\ I_3]$ are totally
unimodular. There are thirteen pairs by
Proposition~\ref{prop:k12-low-excess-simplicial}.

The convex hull of the nonzero signed relation columns is an
origin-symmetric lattice polytope of determinant parameter one.
The complete published classification has five such classes
\cite[\S5.2, Table~2, p.~85]{AverkovSchymura2024}.
Their exact vertex records are the signed lists (TG.4).
The integer-point checks in the accompanying verifier show that each
record contains no additional nonzero lattice points besides its
vertices. Therefore every nonzero signed relation-column class is
exactly one of the columns in one of (TG.4), and every column in that
list occurs. This use of the five-class coverage is a cited primary
input; the complete polytope classification is not claimed rerun.

Write $L=L_0+\mathbb Z(\bar c/3)$. Reduction of the exact presentation
of $L_0$ gives
\[
 \operatorname{row}(B\bmod3)\longrightarrow\mathbb F_3^{13}
      \longrightarrow L_0/3L_0\longrightarrow0.                 \tag{TG.6}
\]
Changing a representative $c$ amounts to adding a row of $B$ modulo
three; replacing its nonzero glue generator also allows multiplication
by $-1$.
Within each nonzero signed Gale class change signs so that its
columns are equal. A transposition of two positions with the same
color fixes $c$, preserves the integral relation lattice and descends
to an automorphism of $L$ exchanging two independent quadratic rays.
It reverses orientation. Thus, for a surviving cell, each such class
has at most three positions, with distinct colors.

On zero Gale columns the coordinate signs may be changed independently.
Among three such positions two colors have the same type, zero or
nonzero. A suitably signed transposition then preserves the glue and
reverses orientation. Thus there are at most two zero columns; with
two their colors must have different types. After sign changes the
possibilities are exactly those recorded above the table.

\emph{The short-support obstruction.}
For any representative of the glue choose its thirteen coefficients
$a_j$ in $\{-1,0,1\}$. Put
$x=\sum a_jv_j/3$ and $t=\#\{j:a_j\ne0\}$.
Since $x\notin L_0$, it is nonzero, and every $x-a_jv_j$ with
$a_j\ne0$ lies outside $L_0$. Normalize the minimum to one.
Completeness of the minimum set implies $Q[x-a_jv_j]>1$.
Direct expansion gives the identity
\[
 \sum_{a_j\ne0}\bigl(Q[x-a_jv_j]-1\bigr)
       =(t-6)Q[x].                                             \tag{TG.7}
\]
The left side is strictly positive. Therefore $t\ge7$.
All twenty-seven row shifts in (TG.6) are tested in the verifier.
A color state with a representative of support at most six cannot
be a complete minimum set. This is an elementary quadratic identity,
not an unverified numerical positivity test.

\emph{All possible stabilizer tests and orbit identifications.}
For a fixed list $\mathcal B=(b_1,\ldots,b_s)$, enumerate every integral
matrix $U$ preserving its signed set. Choose a unit three-column minor;
the images of those three columns must be an ordered signed triple
from the same set with determinant $\pm1$. Testing all such triples
is finite and exhaustive. Write
\[
 Ub_j=\epsilon_j b_{\pi(j)},\qquad\epsilon_j\in\{1,-1\}.
\]
For every $\lambda\in\mathbb F_3^3$ test
\[
 \epsilon_j C(m_j)
     = C(m_{\pi(j)})+\lambda\mathbin{\cdot}b_{\pi(j)}
       \quad\text{for all }j.                                 \tag{TG.8}
\]
If this holds, match each color uniquely by
$c\mapsto\epsilon_jc-\lambda\cdot b_{\pi(j)}$.
The resulting signed coordinate permutation $P$ satisfies
$BP=UB$ and $Pc-c\in\operatorname{row}(B\bmod3)$.
Its transpose preserves the relation lattice, and hence so does $P$.
It descends to $L_0$ and extends to $L$ by (TG.6).
The sign of its ray permutation is exactly its orientation sign.

Conversely, every integral stabilizer of the template permutes its
primitive minimum vectors up to signs. It induces an integral
automorphism of the relation lattice, hence one of the matrices $U$
just enumerated. Its action on the glue line has scalar $1$ or $-1$;
in the latter case compose with global negation, which fixes every
quadratic ray. Signs on zero columns can then be normalized without
changing the orientation permutation. Thus (TG.8) tests every
possible orientation sign. It also gives all equivalences of color
states, together with their row shifts. Distinct published polytope
classes cannot become equivalent in this process.

Here is the full finite enumeration, with masks ranging over
$1,\ldots,7$ and total multiplicity, including zeros, equal to thirteen:
\[
\begin{array}{c|r|r|r|r|r}
\mathcal B_i&\text{states}&\text{odd symmetry}&\text{support}\le6
       &\text{remaining states}&\text{orbits}\\\hline
1&63450&21978&24300&17172&20\\
2&2310&1410&0&900&6\\
3&14&14&0&0&0\\
4&14&14&0&0&0\\
5&0&0&0&0&0\\\hline
\text{total}&65788&23416&24300&18072&26
\end{array}                                                     \tag{TG.9}
\]
The exact verifier exhausts the stated finite loops, verifies every
matrix and color condition with integers, and computes the complete
orbits under (TG.8) and row shifts. Their representatives are precisely
Table~\ref{tab:k12-ternary-templates}. No matrix optimization is used
in this finite classification.

\emph{Exact realization of every retained template.}
The companion data file supplies an integer Gram matrix $H$ and a
positive denominator $d$ for each template, in an explicit basis of
$L$. Its change-of-basis numerator $D$ satisfies
$\det(D/3)=1/3$ and $Dv_j^{\rm ambient}=3v_j^{L_0}$; the glue class
is checked directly modulo three. Thus these are forms on the
specified ambient lattices, not forms only on $L_0$.
For each $H/d$ the verifier constructs and checks its rational
$LDL^t$ decomposition with strictly positive pivots. Descending
enumeration using that identity visits every vector of norm at most
one, with integer square-root bounds and no rounding. Each search
returns exactly the twenty-six signed prescribed vectors. Their
quadratic span has rank thirteen. The searches use between $325$
and $467$ nodes. Complementary minors relative to $L_0$ are all
$0,\pm1$, so the largest basis index in $L$ is exactly three.
These checks prove realization and completeness for all twenty-six
templates. The exhaustive test (TG.8) has no odd stabilizer for any
of them, as asserted.

\emph{Their boundary is zero.}
Each cell is simplicial, so its facets delete a single minimum ray.
A deficient-vector-rank facet is zero in the relative complex.
For a full-rank facet, its twelve minimum vectors still contain a
basis of $L_0$, since every independent ten are such a basis.
It therefore has a rank-two totally unimodular Gale matrix and
the same index-three glue. Heller's bound gives at most three
nonzero signed Gale classes
\cite[\S1, the formula $h(1,m)=m^2+m+1$]{AverkovSchymura2024}.
If it had no odd same-color transposition, each class could contain
at most three positions and the zero class at most two. This accounts
for at most $3\cdot3+2=11$ positions, fewer than twelve.
Thus every full-rank facet has an odd integral stabilizer.
Its orientation coinvariant is zero since two is a unit in $R$.
The cellular incidence sum therefore vanishes.

Finally, minimum-generated cells of maximal index at most three
are already zero or the certified boundary $\sigma$ by
Proposition~\ref{prop:k12-trimodular-generated}; maximal index at most
two vanishes by Proposition~\ref{prop:k12-low-index-orientation}.
These observations prove the final cycle reduction.
\end{proof}

\paragraph{Replayable certificate and precise boundary gate.}
The standard-library verifier
\texttt{k12\_ternary\_glue\_certificate.py} uses the frozen integer
input \texttt{k12\_ternary\_glue\_form\_data.json} and writes its
complete check report as JSON. It checks the five archived records,
the entire finite color enumeration, all twenty-six orbit and
stabilizer conclusions, and every rational minimum certificate.
The only retained classification input is the published completeness
of the five determinant-one polytopes.
To remove this entire contribution to the $K_{12}$ chain gate it is
now sufficient to put these twenty-six explicitly specified cycles
in the incoming boundary image. Such fillings are not supplied by
the classification theorem. Cells of maximal basis index at least
four remain outside its scope.
